\documentclass[12pt,a4paper]{article}

\usepackage[utf8]{inputenc}
\usepackage[T1]{fontenc}
\usepackage[brazil]{babel}
\usepackage{geometry}
\usepackage{hyperref}

\geometry{margin=2.5cm}

\title{Núcleo de Situação de Saúde (NSS)\\Sprint 2 -- Planejamento}
\author{Zadoque Pires de Deus Souza Carneiro}
\date{Agosto de 2026}

\begin{document}

\maketitle

\section{Objetivo da Sprint 2}

A Sprint~2 tem como objetivo principal sair do esqueleto e avançar para uma
primeira versão funcional de ponta a ponta, capaz de:

\begin{itemize}
  \item Executar a pipeline Bronze--Silver--Gold para pelo menos uma doença
        (por exemplo, dengue em 2026), gerando tabelas agregadas compatíveis
        com a API.
  \item Integrar efetivamente a API FastAPI com os arquivos Parquet da camada
        Gold, retornando dados reais para os endpoints de heatmap.
  \item Evoluir o front-end para uma visualização mais próxima do mapa de
        calor desejado, iniciando o uso de malhas geográficas oficiais.
\end{itemize}

\section{Atividades planejadas}

\subsection{Integração da pipeline Bronze--Silver--Gold}

\begin{itemize}
  \item \textbf{Bronze}: incorporar os scripts de ingestão existentes (via
        PySUS) ao repositório, organizando-os em um pacote Python dentro do
        diretório \texttt{back-end/} (ou em uma pasta dedicada
        \texttt{data-pipeline/}).
  \item \textbf{Silver}: adaptar os scripts genéricos de Silver para seguir a
        convenção de colunas definida na Sprint~1, incluindo normalização de
        códigos de município/UF e deduplicação de registros.
  \item \textbf{Gold}: implementar um agregador que, a partir da Silver,
        produza tabelas com colunas \texttt{disease\_code}, \texttt{year},
        \texttt{month}, \texttt{cd\_uf}, \texttt{nm\_uf}, \texttt{cd\_mun},
        \texttt{nm\_mun} e \texttt{cases\_total}, salvando em Parquet na
        estrutura \texttt{gold/sinan/...}.
  \item Automatizar, quando possível, a execução dessas etapas (por exemplo,
        via scripts \texttt{Makefile}, \texttt{tox} ou mesmo \texttt{Docker
        Compose}).
\end{itemize}

\subsection{Integração da API com os dados Gold}

\begin{itemize}
  \item Ajustar as funções de leitura da camada Gold em
        \texttt{app/main.py} para utilizar os arquivos produzidos pela
        pipeline.
  \item Garantir que os endpoints \texttt{/heatmap/uf} e
        \texttt{/heatmap/municipios} utilizem efetivamente as colunas
        \texttt{year}, \texttt{month}, \texttt{cd\_uf}, \texttt{cd\_mun} e
        \texttt{cases\_total} para filtrar e agregar os dados.
  \item Adicionar testes básicos (unitários ou de integração) para validar o
        comportamento da API quando existem dados e quando a tabela está
        vazia.
  \item Documentar, no \texttt{README.md}, exemplos de chamadas reais à API
        com dados de dengue.
\end{itemize}

\subsection{Evolução do front-end}

\begin{itemize}
  \item Substituir (ou complementar) o protótipo estático atual por uma
        aplicação web mais estruturada, preferencialmente utilizando React ou
        outra biblioteca com a qual o time tenha familiaridade.
  \item Integrar uma biblioteca de mapas (por exemplo, Leaflet ou MapLibre)
        para exibir um mapa do Brasil com as UFs.
  \item Realizar o \emph{join} entre os dados retornados pelo endpoint
        \texttt{/heatmap/uf} e a malha geográfica, aplicando uma escala de
        cores para representar a intensidade dos casos (mapa de calor).
  \item Implementar interação básica: \emph{hover} para mostrar o número de
        casos de uma UF e \emph{click} para navegar para o nível de
        municípios daquela UF (mesmo que o mapa em nível municipal ainda seja
        simplificado nesta sprint).
\end{itemize}

\subsection{Malhas geográficas e códigos IBGE}

\begin{itemize}
  \item Baixar as malhas territoriais oficiais do IBGE (UFs e municípios) no
        formato apropriado (Shapefile ou GeoJSON).
  \item Garantir a consistência entre os códigos de município/UF usados na
        pipeline Silver/Gold e os códigos presentes nas malhas.
  \item Preparar arquivos GeoJSON simplificados (por exemplo, com redução de
        resolução) para uso eficiente no front-end.
\end{itemize}

\section{Critérios de aceitação}

A Sprint~2 será considerada bem-sucedida se:

\begin{itemize}
  \item For possível executar a pipeline completa para uma doença-alvo e
        visualizar, via API, os agregados por UF.
  \item O protótipo de mapa de calor por UF estiver funcional, consumindo
        dados reais da API.
  \item A documentação (\texttt{README.md} e os arquivos de sprint) estiver
        atualizada com instruções claras de execução.
\end{itemize}

\section{Riscos e observações}

Alguns riscos e pontos de atenção identificados:

\begin{itemize}
  \item Complexidade na limpeza dos dados Bronze (muitas colunas com valores
        nulos, diferentes versões de layout do SINAN, etc.).
  \item Tamanho dos arquivos de malha municipal do IBGE, que pode impactar o
        desempenho do front-end caso não haja simplificação.
  \item Necessidade de padronização dos códigos de município/UF ao longo de
        toda a pipeline para evitar erros de \emph{join} entre dados e
        geometria.
\end{itemize}

Apesar desses riscos, a Sprint~2 foca em resultados incrementais, permitindo
ajustes e refinos nas sprints seguintes.

\end{document}
